\documentclass[11pt,a4paper]{article}
\usepackage[margin=2.4cm]{geometry}
\usepackage{hcqa-style}

\hypersetup{
  pdftitle={hi\_cube\_qa -- A Quality-Assessment Toolkit for MeerKAT HI Spectral-Line Cubes},
  pdfauthor={GAMA G23 MeerKAT HI Survey}
}

\begin{document}

% --------------------------- title block ----------------------------------
\begin{center}
  {\hcqalogo[1.7]}\\[14pt]
  {\LARGE\bfseries\ttfamily hi\_cube\_qa}\\[6pt]
  {\large A Quality-Assessment Toolkit for MeerKAT \textsc{Hi}\ Spectral-Line Cubes}\\[10pt]
  {\normalsize GAMA~G23 MeerKAT \textsc{Hi}\ Survey --- Data-Reduction QA}\\[2pt]
  {\normalsize Reference Manual \quad\textbullet\quad \today}
\end{center}
\noindent\rule{\textwidth}{0.6pt}
\vspace{6pt}

\begin{abstract}
\noindent
\code{hi\_cube\_qa} is a Python toolkit that turns the spectral-line cubes
produced by the \textsc{CARACal} pipeline \citep{jozsa2020} into a single,
self-contained, interactive HTML quality-assessment (QA) report \emph{per
observation session}. It auto-discovers every cube under a session's directory
tree, measures the per-channel noise, the synthesised-beam behaviour, and the
spectral response of each imaging product (\code{dirty}, \code{image},
\code{residual}, \dots), compares the measured noise against the
radiometer-equation expectation, and flags channels and cubes that fall outside
expert-defined tolerances. This document describes the scientific motivation
behind every metric, the exact quantities the code computes, and how to
configure and run the toolkit. A companion \emph{Quickstart} covers the common
workflow in two pages.
\end{abstract}

\tableofcontents
\bigskip

% ===========================================================================
\section{Purpose and scope}
% ===========================================================================
Radio-interferometric \textsc{Hi}\ cubes are large ($10^2$--$10^3$ channels,
each a $\sim$\,$10^3\times10^3$-pixel image) and their scientific usefulness is
limited by a handful of measurable quantities: how close the achieved noise is
to the theoretical limit, how that noise varies with frequency (revealing
residual radio-frequency interference, RFI, and calibration errors), how stable
the synthesised beam is across the band, and how much data the flagging stages
removed. Inspecting these by eye, cube by cube, is slow and error-prone.

\code{hi\_cube\_qa} automates that inspection. It is built for the GAMA~G23
MeerKAT \citep{jonas2016} \textsc{Hi}\ programme, in which each observation
\emph{session} (one correlator dump, identified by a $10$-digit block ID such as
\code{1610268336\_sdp\_l0}) is imaged into several \emph{pointings}, each
producing one or more spectral \emph{cubes} in several imaging \emph{products}.
The toolkit produces one HTML report per session so that the quality of an
entire night's imaging can be reviewed in a single document.

% ===========================================================================
\section{Data model and pipeline architecture}\label{sec:datamodel}
% ===========================================================================
\subsection{\texorpdfstring{The session $\to$ pointing $\to$ unit $\to$ product
hierarchy}{The session, pointing, unit and product hierarchy}}
\textsc{CARACal} writes its \textsc{Hi}\ output into a fixed directory layout:
\begin{lstlisting}
<obsid>_sdp_l0/caracal/
    output_hi_<pXX>/                     # one imaging run per pointing
        cubes/
            cube_1/ G23_<obsid>_<pXX>_HI_1.<product>.fits
            cube_2/ G23_<obsid>_<pXX>_HI.<product>.fits
    output_hi_1gc/diagnostic_plots/      # flagging-summary JSONs
\end{lstlisting}
The toolkit mirrors this structure with four nested concepts:
\begin{description}[leftmargin=1.6em,style=nextline]
  \item[Session] One \code{caracal/} root; produces exactly one HTML report.
  \item[Pointing] One \code{output\_hi\_pXX} imaging run (e.g.\ \code{p63}).
  \item[Unit] One \code{(pointing, cube\_N)} pair. \emph{All imaging products
        inside a single \code{cube\_N} share the same spatial WCS and channel
        grid}, so they are legitimately overlaid on common axes. Different
        \code{cube\_N} (different spectral windows) are kept as separate units.
  \item[Product] One FITS file: \code{dirty}, \code{image}, \code{residual},
        and optionally \code{model}, \code{psf}.
\end{description}
This ``unit'' abstraction is the key design decision: it guarantees that every
overlaid comparison plot mixes only cubes that genuinely share an $x$-axis,
while still letting a single report span many pointings.

\subsection{Auto-discovery (\code{discovery.py})}
Given a session root, \code{discover\_session()} globs
\filepath{output_hi_p*/cubes/cube_*/*.fits}, parses the imaging product from the
filename suffix (the alphabetic token before \code{.fits}), and groups the files
into units. The user never lists individual paths; they choose only \emph{which
products}, and optionally which \emph{pointings} or \emph{cubes}
(\code{cube\_N}), to keep (Section~\ref{sec:config}). By default every
\code{cube\_N} found in each pointing becomes its own unit. Products are ordered
consistently
(\code{dirty}, \code{image}, \code{residual}, \code{model}, \code{psf}) so that
plot legends and colours are stable across the whole report.

\subsection{Spectral-axis and WCS handling (\code{io\_utils.py})}
FITS cubes from \textsc{CARACal}/WSClean \citep{offringa2014} may store the
spectral axis as frequency (\code{FREQ}, Hz) or as velocity (\code{VELO},
\code{VRAD}, \code{VOPT}, in \si{m/s} or \si{km/s}). \code{load\_cube()} locates
the spectral axis, reads its world-coordinate keywords, and constructs a matched
frequency/velocity pair. The radio velocity convention is used throughout,
\begin{equation}
  v_{\mathrm{radio}} = c\,\frac{f_0 - f}{f_0},
  \qquad
  f = f_0\!\left(1 - \frac{v_{\mathrm{radio}}}{c}\right),
  \label{eq:radio}
\end{equation}
with the \textsc{Hi}\ rest frequency $f_0 = \SI{1.420405752}{GHz}$ (or
\code{RESTFRQ} from the header) and $c = \SI{2.99792458e5}{km/s}$. Optical
velocities are converted via $f = f_0/(1+v_{\mathrm{opt}}/c)$ and then expressed
as radio velocity for internal consistency. The channel width in hertz,
\begin{equation}
  \Delta\nu =
  \begin{cases}
    |\mathrm{CDELT}_{\mathrm{spec}}| & \text{(\code{FREQ} axis),}\\[2pt]
    f_0\,|\Delta v|/c & \text{(velocity axis),}
  \end{cases}
\end{equation}
is read directly from each cube's header and \emph{overrides} the configured
default, so the theoretical-noise prediction (Section~\ref{sec:radiometer})
always uses the true spectral resolution of the cube being assessed.

% ===========================================================================
\section{Noise metrics}\label{sec:noise}
% ===========================================================================
Noise is the single most important QA quantity for an \textsc{Hi}\ cube: it sets
the column-density sensitivity and directly exposes calibration and RFI
problems. The toolkit measures noise robustly, channel by channel, and compares
it against first-principles expectations.

\subsection{Robust per-channel RMS}\label{sec:madrms}
For each spectral plane the noise is estimated with the median absolute
deviation (MAD) rather than the ordinary standard deviation:
\begin{equation}
  \sigma_{\mathrm{MAD}} = 1.4826 \times \operatorname{median}_i
  \bigl|\,x_i - \operatorname{median}(x)\,\bigr|.
  \label{eq:mad}
\end{equation}
\paragraph{Why MAD.} A spectral plane contains not only noise but also real
\textsc{Hi}\ emission, continuum sources, deconvolution residuals and RFI
spikes. The ordinary standard deviation is dominated by these outliers and
\emph{overestimates} the noise. The MAD is a high-breakdown-point estimator
(up to $50\%$ of the data can be contaminated before it fails), so it returns
the noise of the bulk of ``empty'' pixels even when bright sources are present.
The factor $1.4826 = 1/\Phi^{-1}(3/4)$ rescales the MAD so that, for Gaussian
noise, $\sigma_{\mathrm{MAD}}$ equals the true standard deviation; the reported
quantity is therefore directly comparable to the radiometer prediction.

\subsection{Peak per-channel amplitude}
As a complementary tracer the code records the per-channel maximum absolute
pixel value, $\max_{x,y}|T(x,y)|$. This is a proxy for the worst-case artefact
or unflagged RFI spike in each channel and, divided by
$\sigma_{\mathrm{MAD}}$, yields the dynamic range / peak signal-to-noise ratio
used to pick diagnostic channel maps (Section~\ref{sec:thumbs}).

\subsection{Boresight (best-case) noise}\label{sec:boresight}
The full-image MAD RMS is a \emph{primary-beam-weighted average} over a field
whose sensitivity falls towards the edges. To obtain the \emph{best-case} noise
for comparison with theory, \code{boresight\_rms()} measures the RMS in a square
box (default $100$ pixels) centred on the pointing centre (\code{CRPIX}), where
the primary-beam response is highest. Source pixels are rejected with iterative
MAD-based $\sigma$-clipping:
\begin{enumerate}[noitemsep,topsep=2pt]
  \item compute the median and $\sigma_{\mathrm{MAD}}$ of the surviving pixels;
  \item reject pixels deviating by more than $\kappa\,\sigma_{\mathrm{MAD}}$
        (default $\kappa = 3$);
  \item repeat (default $3$ iterations).
\end{enumerate}
The MAD-RMS of the surviving pixels is reported per channel (the same robust
estimator as the full-image RMS, for consistency). Because it is measured where
the telescope is most sensitive and after outlier rejection, it is the fairest
measured quantity to set against the radiometer limit. By default the
\code{residual} product (if present) is used as the noise reference here, as it
is the cleanest noise tracer (Section~\ref{sec:noiseref}). The box is read
directly from the memory-mapped cube (Section~\ref{sec:performance}), so this
costs no extra full-cube pass.

\subsection{The radiometer-equation expectation}\label{sec:radiometer}
The theoretical thermal noise of a single channel of an interferometric image
\citep{tms2017,nraooss} is computed by \code{theoretical\_noise()}. The
system-equivalent flux density (SEFD) is taken directly from the config (the default is
\SI{425}{Jy} per antenna, the value SARAO assumes in its L-band sensitivity calculator,
\citealp{saraosens}) or derived from the system temperature and aperture efficiency,
\begin{equation}
  \mathrm{SEFD} = \frac{2 k_{\mathrm B}\,T_{\mathrm{sys}}}{\eta\,A_{\mathrm{geo}}}\times 10^{26}\ \si{Jy},
  \qquad
  A_{\mathrm{geo}} = \pi\!\left(\frac{d}{2}\right)^{\!2},
  \label{eq:sefd}
\end{equation}
with the MeerKAT dish diameter $d = \SI{13.5}{m}$, Boltzmann constant
$k_{\mathrm B}$, and the factor $10^{26}$ converting
\si{W.m^{-2}.Hz^{-1}}\ to \si{Jy}. The per-channel point-source noise is then
\begin{equation}
  \sigma_{\mathrm{theory}} =
  \frac{n_{\mathrm{fac}}\;\mathrm{SEFD}}
       {\sqrt{\,n_{\mathrm{pair}}\;\tau\;\Delta\nu\;n_{\mathrm{pol}}\,}},
  \qquad
  n_{\mathrm{pair}} = N_{\mathrm{ant}}\,(N_{\mathrm{ant}}-1),
  \label{eq:radiometer}
\end{equation}
where $n_{\mathrm{pair}}$ is the number of ordered antenna pairs (twice the number of
baselines), $\tau$ the on-source
integration time (seconds), $\Delta\nu$ the channel width (hertz),
$n_{\mathrm{pol}} = 2$ the number of polarisations averaged, and
$n_{\mathrm{fac}}$ a dimensionless weighting penalty (below). Equation
(\ref{eq:radiometer}) is the natural-weighting thermal-noise expression of the
NRAO VLA Observational Status Summary \citep{nraooss}, with its correlator efficiency
taken as 1 (no MeerKAT value is available here). It is
retained as a band-centre scalar for logging and for the band-averaged grade,
but the comparison drawn against the per-channel RMS uses the
\emph{frequency- and flagging-aware} form of the next two subsections.

\paragraph{Weighting penalty $n_{\mathrm{fac}}$.} Imaging weights trade angular
resolution against sensitivity \citep{briggs1995}: natural weighting gives the
lowest noise, uniform the best resolution, Briggs robust interpolates. The
penalty $n_{\mathrm{fac}}$ scales the noise above natural. It can be set
explicitly, or \code{nfac: auto} reads the Briggs \code{robust} value from the
header (Section~\ref{sec:provenance}) and maps it to an approximate penalty, so
the prediction matches the weighting actually used.

\paragraph{Number of antennas $N_{\mathrm{ant}}$.} Since $n_{\mathrm{pair}}$
enters Eq.~(\ref{eq:radiometer}) as $N_{\mathrm{ant}}(N_{\mathrm{ant}}-1)$, a
fixed assumed array size is a poor approximation across a multi-session
campaign: individual sessions can differ in how many antennas were available
(e.g.\ some offline for maintenance), and the theory should reflect the array
that was actually used. With \code{radiometer.n\_ant: auto} (the default), the
real per-session antenna count is read from the CARACal \code{obsinfo}
summary's \code{ANT} field once per session; sessions without this information
fall back to \code{n\_ant\_fallback} (default $60$). A session with $10$ fewer
antennas than assumed changes $\sigma_{\mathrm{theory}}$ by a factor
$\sqrt{n_{\mathrm{pair}}^{\mathrm{assumed}}/n_{\mathrm{pair}}^{\mathrm{real}}}$ --
e.g.\ $60\to50$ antennas raises the predicted noise by $\sim\!20\%$ -- so
getting this right matters for the RMS/theory comparison to be meaningful.

\code{ANT} is keyed by MS ANTENNA-table \emph{column} name (\code{NAME},
\code{STATION}, \code{OFFSET}, \code{POSITION}, \code{DISH\_DIAMETER},
\code{MOUNT}, \code{TYPE}, \code{FLAG\_ROW}), each mapping to a list with one
entry \emph{per antenna} -- confirmed directly on real G23 data, where the
dict has exactly these $8$ top-level keys (the fixed column count) while, e.g.,
its \code{NAME} column held $58$ entries (the real antenna count). An earlier
version of this reader took \code{len(ANT)}, silently returning the column
count ($8$) rather than the antenna count, and was caught only once a live run
showed an implausible $N_{\mathrm{ant}}=8$. The reader now takes the length of
whichever per-antenna column is present, in a fixed preference order.

\subsection{Frequency-dependent SEFD}\label{sec:sefd}
A single $T_{\mathrm{sys}}$ across the whole band is a poor approximation: the
MeerKAT L-band SEFD has a broad minimum near \SI{1.3}{GHz} and rises towards both
edges (and around the satellite-navigation allocations). Using a flat value
makes the RMS-to-theory ratio spuriously large at the band edges even for good
data. \code{sefd\_spectrum()} therefore evaluates an SEFD \emph{curve}
$\mathrm{SEFD}(\nu)$, normalised so its band-centre value equals the
$T_{\mathrm{sys}}/\eta$ SEFD of Eq.~(\ref{eq:sefd}) --- i.e.\ it adopts the
\emph{shape} of an illustrative L-band curve while keeping the user's
$T_{\mathrm{sys}}$ as the absolute scale. The curve is configurable
(\code{radiometer.sefd\_curve}: \code{default}, \code{flat}, or a list of
$[\nu_{\mathrm{MHz}}, \mathrm{SEFD}_{\mathrm{Jy}}]$ points).

\paragraph{Caveat.} The curve is used only when \code{radiometer.sefd} is null; with the
default flat \SI{425}{Jy} it is ignored. The built-in default curve is \emph{illustrative} ---
it captures the band shape but is not the official SARAO calibration. For
quantitative work, supply the measured SEFD$(\nu)$ for the observations being assessed via
\code{sefd\_curve}, or set it to \code{flat} to fall back to the single
$T_{\mathrm{sys}}$.

\subsection{Flagging-aware effective time}\label{sec:flagaware}
RFI flagging removes visibilities, reducing the \emph{effective} on-source time
and therefore raising the expected noise. The per-channel theory used for
classification folds this in,
\begin{equation}
  \sigma_{\mathrm{theory}}(\nu) =
  \frac{n_{\mathrm{fac}}\,\mathrm{SEFD}(\nu)}
       {\sqrt{\,n_{\mathrm{pair}}\;\tau_{\mathrm{eff}}(\nu)\;\Delta\nu\;n_{\mathrm{pol}}\,}},
  \qquad
  \tau_{\mathrm{eff}}(\nu) = \tau\,\bigl(1 - f_{\mathrm{flag}}(\nu)\bigr).
  \label{eq:theoryflag}
\end{equation}
A heavily-flagged sub-band is then \emph{expected} to be noisier, so the
RMS/theory comparison no longer penalises channels for flagging that already
happened. The noise-overview plot shows both the flat and the flagging-reduced
theory curves so the effect is explicit.

\paragraph{Per-pointing matching.} The CARACal flag reports are written
\emph{per pointing}, with the pointing id encoded in the filename (e.g.\ a
\code{p20} \code{target\_summary} report). With \code{flag\_fraction: auto} the
toolkit reads the overall flagged fraction from each pointing's
\code{target\_summary} report. The CARACal report is a \code{Flag stats} list of
categories (\code{scans}, \code{antennas}, \code{fields}, \code{corrs}), each
partitioning the same data into entries carrying \code{sum} (flagged) and
\code{counts} (total); the overall fraction is
$\sum\code{sum}/\sum\code{counts}$ over any one category. That fraction is
applied to \emph{that pointing's} cubes; pointings without a report fall back to
the session median. A fixed number or \code{null} disables the per-pointing behaviour, and
each unit's fraction is shown in its provenance strip.

\paragraph{Caveat.} The flag fractions are still \emph{scalar} per pointing, not
a per-channel flag spectrum, so the reduction in $\tau_{\mathrm{eff}}$ is
uniform across the band within a pointing. A genuine $f_{\mathrm{flag}}(\nu)$
would require the flag cube from the measurement set, which is not among these
products.

\paragraph{Automatic discovery.} Flag reports are found per session by globbing
\code{flagreport\_glob} relative to the session root; the default matches
\code{output\_hi\_1gc/diagnostic\_plots/}. If that finds nothing (a session
with a different layout), discovery falls back to progressively broader globs
(\code{output\_hi\_*gc/}, any \code{output\_hi\_*/diagnostic\_plots/}, then a
recursive search), so no per-session configuration is needed. The matched
pattern is logged, and a session with no flags found degrades gracefully (the
theory is simply not flag-reduced).

\subsection{Real per-channel data loss: investigated, abandoned}\label{sec:nvislog}
The scalar flagged fraction of Section~\ref{sec:flagaware} is a coarse proxy:
it reduces $\tau_{\mathrm{eff}}$ uniformly across the whole band, whereas real
RFI is concentrated in specific sub-bands. Since CARACal's per-pointing
imaging step invokes WSClean \citep{offringa2014} once per cube, covering the
full channel range, and WSClean logs a line
\begin{quote}\ttfamily Gridded visibility count: \textless N\textgreater\end{quote}
each time it (re-)grids a channel, this looked like a route to a genuine
per-channel data-loss quantity -- exactly what the CHILES survey uses for its
own noise validation \citep{luber2025} -- read directly from the pipeline's own
logs rather than a separate measurement.

This was implemented and then withdrawn after testing against real G23 data.
WSClean's message is generic and carries no channel index, so a parser must
infer channel boundaries from where the logged value \emph{changes}. That
assumption holds for a short sample, but breaks down over a full band: a
contiguous run of fully-flagged channels can log the identical value for
hundreds of consecutive lines, with nothing in the log to distinguish "one
empty channel" from "many in a row." On a real 4783-channel cube, the typical
run length was 4--5 lines, but one run reached $1442$ -- two orders of
magnitude larger than anything else -- silently collapsing what was almost
certainly dozens to hundreds of distinct channels into one. This is not a
parsing bug fixable with a smarter heuristic; the information needed to
disambiguate is simply absent from the log text. The flagged-fraction theory
of Section~\ref{sec:flagaware} and the kurtosis diagnostic
(Section~\ref{sec:kurtosis}) remain the validated tools for this purpose.

\subsection{Per-channel quality classification}\label{sec:classify}
Each channel is labelled \textsf{Good}, \textsf{Warn} or \textsf{Bad} by
\code{classify\_channels()} from the ratio of its measured RMS to a reference.
If a valid theoretical noise is available it is the reference; otherwise the
band-median RMS is used (so the tool still works without radiometer inputs).
With warn/bad factors $w = 1.5$ and $b = 3.0$ (configurable),
\begin{equation}
  \text{channel}\ i \rightarrow
  \begin{cases}
    \textsf{Bad}  & \sigma_i\ \text{non-finite, or}\ \sigma_i/\sigma_{\mathrm{ref}} > b,\\
    \textsf{Warn} & w < \sigma_i/\sigma_{\mathrm{ref}} \le b,\\
    \textsf{Good} & \sigma_i/\sigma_{\mathrm{ref}} \le w.
  \end{cases}
\end{equation}
A non-finite RMS marks a fully flagged or empty plane as \textsf{Bad}. The
band-level summary (\code{channel\_summary()}) aggregates these into counts and
percentages and assigns an overall cube grade:
\begin{equation}
  \text{cube} \rightarrow
  \begin{cases}
    \textsf{Bad}  & f_{\mathrm{bad}} \ge 15\%\ \text{or}\ R > 3,\\
    \textsf{Warn} & f_{\mathrm{bad}} \ge 5\%\ \text{or}\ f_{\mathrm{warn}} \ge 20\%\ \text{or}\ R > 1.5,\\
    \textsf{Good} & \text{otherwise,}
  \end{cases}
\end{equation}
where $f_{\mathrm{bad}}$, $f_{\mathrm{warn}}$ are the bad/warn channel fractions
and $R$ is the band-averaged RMS-to-theory ratio (the median over the
\textsf{Good} channels, or over all finite channels if none are Good, divided by
$\sigma_{\mathrm{theory}}$).

\paragraph{Interpreting $R\sim1.5$.} A ratio somewhat above unity is normal, not
a fault: real MeerKAT \textsc{Hi}\ imaging typically achieves $\sim1.2$--$1.5$
times the thermal prediction because of residual calibration errors, confusion,
and imaging losses not captured by the idealised radiometer equation. A whole
band sitting at $R\sim1.5$ (many \textsf{Warn} channels) is therefore expected
for good data; it is a \emph{large} excess or a rising/structured RMS spectrum
that signals a real problem. If \emph{every} channel is \textsf{Warn}/\textsf{Bad}
the median RMS falls back to all finite channels (rather than reporting nothing),
so a uniformly noisy but valid cube still reports its true noise.

\subsection{Why the residual is the noise reference}\label{sec:noiseref}
For the noise-overview comparison, the boresight, the 2-D maps and the
band-averaged grade the toolkit prefers, in order, \code{residual} $\to$
\code{dirty} $\to$ \code{image}.
The \code{residual} cube is the dirty image minus the restored clean model: it
contains the thermal noise plus un-deconvolved sidelobes but \emph{not} the
clean components, making it the closest available estimate of the true noise
floor. The \code{image} cube, by contrast, contains restored source flux that
can bias a naive RMS upward; it is still measured and displayed, but is not used
as the noise yardstick.

\subsection{Worked example: a deconvolution failure is visible only in the cleaned products}
\label{sec:divergence_example}
Comparing the products channel by channel separates problems in the data from problems
introduced by deconvolution. In the October 2026 G23 processing, isolated channels of the
second-pass (\textsc{SoFiA}-masked) \code{cube\_2} of several pointings had a \code{residual}
RMS of $10^2$--$10^3$ times the normal level (for example $5.4$\,Jy against $0.97$\,mJy in one
channel at $1362.0$\,MHz), while the \code{dirty} cube, which contains no deconvolution, and
the first-pass \code{cube\_1} were normal in the same channels. A channel that is high in all
products is a data problem (RFI, calibration); a channel that is high only in the cleaned
products is a deconvolution failure. Here the cause was WSClean's automatic multiscale list
reaching a scale as large as the image, which oscillates and grows inside the mask; capping the
scales fixed it (Sect.~\ref{sec:noiseref} explains why the \code{residual} is the right cube to
look at). This is why the default products include \code{residual}: with \code{dirty} and
\code{image} alone the failure is easy to miss. The failure is also not always flagged by the
deconvolver's own warnings, so the per-channel RMS, not the log, should decide whether a cube
is usable.

% ===========================================================================
\section{Synthesised-beam diagnostics}\label{sec:beam}
% ===========================================================================
The synthesised beam (point-spread function) determines the angular resolution
and the conversion between surface brightness (\jybeam) and flux density (Jy).
For wide-band \textsc{Hi}\ imaging the beam changes with frequency, and a beam
that varies erratically across the band signals deconvolution or $uv$-coverage
problems. \code{beam\_stats()} reads the beam from the FITS header, supporting
three conventions in priority order:
\begin{enumerate}[noitemsep,topsep=2pt]
  \item single global \code{BMAJ}/\code{BMIN}/\code{BPA} keywords, broadcast to
        all channels;
  \item per-channel WSClean keywords \citep{offringa2014} (\code{BMAJ0000}\dots,
        zero-padded $0$-indexed; or \code{BMAJ1}\dots, $1$-indexed, with
        $8$-character FITS keyword truncation handled, e.g.\ \code{BMA1});
  \item a \code{BEAMTABLE} binary extension (reserved; currently falls through
        to a warning).
\end{enumerate}
From the per-channel major/minor axes ($\theta_{\mathrm{maj}}$,
$\theta_{\mathrm{min}}$, arcsec) and position angle, \code{beam\_\allowbreak diagnostics()}
reports the band medians and scatter, the beam solid angle
\begin{equation}
  \Omega_{\mathrm{beam}} = \frac{\pi}{4\ln 2}\,\theta_{\mathrm{maj}}\,\theta_{\mathrm{min}},
  \label{eq:beamarea}
\end{equation}
(the integral of a 2-D Gaussian of these FWHMs), and the axial ratio
$\theta_{\mathrm{maj}}/\theta_{\mathrm{min}}$. A channel-to-channel standard
deviation of $\theta_{\mathrm{maj}}$ exceeding $\SI{0.5}{arcsec}$ (configurable)
raises a beam-stability warning, since a large scatter undermines a common
restoring beam and complicates moment analysis.

% ===========================================================================
\section{Non-Gaussianity: per-channel kurtosis}\label{sec:kurtosis}
% ===========================================================================
The RMS-based checks of Section~\ref{sec:noise} test only the \emph{second
moment} of the pixel distribution: a channel can have a perfectly acceptable
RMS while its distribution is heavily non-Gaussian, e.g.\ a small number of
extreme outlier pixels from residual RFI or incomplete continuum subtraction.
Following the validation approach of the CHILES survey \citep{luber2025}, the
toolkit therefore also computes the \emph{fourth moment} -- the Pearson
kurtosis -- per channel, as an independent shape diagnostic:
\begin{equation}
  \kappa = \frac{\left\langle (x-\mu)^4\right\rangle}{\left\langle (x-\mu)^2\right\rangle^2},
  \label{eq:kurtosis}
\end{equation}
using the \emph{ordinary} mean $\mu$ and variance (not the robust MAD
statistics used elsewhere in this toolkit) -- kurtosis exists specifically to
be sensitive to outliers and tails, so a robust, outlier-resistant estimator
would defeat its purpose. A Gaussian distribution has $\kappa = 3$ in this
convention (matching \citealt{luber2025}); values above $3$ (leptokurtic)
indicate heavier tails than Gaussian, values below (platykurtic) indicate
fewer extreme values than expected.

\subsection{Significance and the beam-correlation correction}
Under normality, the sampling variance of $\kappa$ is approximately $24/N$ for
$N$ \emph{independent} samples. Neighbouring pixels within a synthesised beam
are correlated, however, so the number of independent samples is not the raw
pixel count but the number of independent \emph{resolution elements},
$N_{\mathrm{eff}} = N_{\mathrm{pix}}/n_{\mathrm{pb}}$ (pixels per beam,
Section~\ref{sec:bandspec}). The toolkit therefore uses
\begin{equation}
  \mathrm{SE}(\kappa) = \sqrt{\frac{24}{N_{\mathrm{eff}}}},
  \label{eq:kurtosisse}
\end{equation}
and classifies a channel as \textsf{Warn} or \textsf{Bad} when $|\kappa - 3|$
exceeds $5$ or $10$ standard errors, respectively
(\code{thresholds.kurtosis\_\allowbreak warn\_sigma},
\code{kurtosis\_\allowbreak bad\_sigma}).
\emph{Caveat:} \citet{luber2025} state an
expected kurtosis variation of $0.004$ for their data, which by
Eq.~(\ref{eq:kurtosisse}) implies $N_{\mathrm{eff}} \approx 1.5\times10^6$ --
suspiciously close to their raw $1200\times1200$ pixel count. This suggests
their published uncertainty may not itself carry the beam-correlation
correction applied here; the toolkit's threshold is therefore intentionally
more conservative (wider SE, so fewer false positives) than a literal reading
of that reference value would give.

\subsection{Cross-referencing with the RMS check}
The most useful result is not the kurtosis curve alone but its
\emph{combination} with the RMS/theory check: for every product, the toolkit
counts channels that are \textsf{Bad} on kurtosis but \textsf{Good} on RMS
(non-Gaussian contamination the RMS check cannot see), channels \textsf{Bad} on
both (genuinely catastrophic channels), and channels \textsf{Bad} on RMS but
still Gaussian (e.g.\ a fully flagged plane). This mirrors
\citeauthor{luber2025}'s own central finding -- that their RMS ratio flagged a
narrow RFI-dominated sub-band while kurtosis revealed mild non-Gaussianity
across the \emph{entire} band -- two different failure modes that only appear
together in the cross-reference. The counts are logged per unit and appear as
the \textsf{Gaussian noise (kurtosis)} checklist criterion's note field.

% ===========================================================================
\section{Spectral products}\label{sec:spectral}
% ===========================================================================
\subsection{Spectral profile: peak and mean}
\code{fig\_spectral\_profile()} plots, per channel, the peak
($\max_{x,y}|T|$) and mean flux. The peak trace highlights bright emission or
artefacts; the mean trace (typically at the \si{\micro Jy}\ level) is sensitive
to large-scale additive errors such as continuum-subtraction residuals or
baseline ripples.

\subsection{Integrated band spectrum}\label{sec:bandspec}
The spatially integrated flux-density spectrum converts the cube from surface
brightness to flux density. With pixel solid angle
$\Omega_{\mathrm{pix}} = |\mathrm{CDELT}_1|\,|\mathrm{CDELT}_2|$ (in steradian)
and beam solid angle $\Omega_{\mathrm{beam}}$ from Eq.~(\ref{eq:beamarea}), the
number of pixels per beam is $n_{\mathrm{pb}} = \Omega_{\mathrm{beam}}/\Omega_{\mathrm{pix}}$, and
\begin{equation}
  S_{\mathrm{int}}(\nu) = \frac{1}{n_{\mathrm{pb}}}\sum_{x,y} T(x,y,\nu)
  \quad[\si{Jy}].
  \label{eq:sint}
\end{equation}
Dividing the summed \jybeam\ by $n_{\mathrm{pb}}$ removes the multiple-counting
of flux that is spread over many pixels within one beam, yielding a true flux
density. The uncertainty envelope on the integrated spectrum is propagated from
the per-channel noise over the typical number of finite pixels $N_{\mathrm{pix}}$,
\begin{equation}
  \sigma_{\mathrm{int}} = \sigma_{\mathrm{MAD}}\sqrt{\frac{N_{\mathrm{pix}}}{n_{\mathrm{pb}}}},
\end{equation}
and $\pm1\sigma$ and $\pm3\sigma$ bands are shaded. The plot carries a dual
abscissa (radio velocity and frequency, related by Eq.~\ref{eq:radio}) and can
mark a target redshift $z$. The marker is placed at the correct \emph{radio}
velocity of the line, $v = cz/(1+z)$ (with $f = f_0/(1+z)$); using the naive
$v = cz$ would misplace it by a factor $(1+z)$ --- a few hundred \si{km/s} at the
$G23$ redshifts.

\subsection{Diagnostic channel maps}\label{sec:thumbs}
Rather than rendering every channel, \code{thumbnails.py} shows a small set of
\emph{diagnostic} channel images chosen on two physical axes
(\code{thumbnails.selectors}):
\begin{itemize}[noitemsep,topsep=2pt]
  \item \textbf{by RMS} --- the cleanest (minimum), typical (median) and
        noisiest (maximum) channels, exposing the best and worst noise
        behaviour and any catastrophic planes;
  \item \textbf{by SNR} --- the strongest (peak), typical (median) and weakest
        (minimum) signal-to-noise channels, where
        $\mathrm{SNR} = \max_{x,y}|T|/\sigma_{\mathrm{MAD}}$, letting the user
        eyeball the most significant detection in the cube.
\end{itemize}
Channels selected by more than one criterion are shown once, annotated with all
the reasons they were chosen. Each map is displayed with a fixed
$[-3\sigma, +5\sigma]$ stretch (from the band-median RMS), a beam ellipse, and a
quality badge, so panels are directly comparable. This keeps reports small while
still surfacing exactly the channels a reviewer needs to see.

% ===========================================================================
\section{HI column-density sensitivity}\label{sec:nhi}
% ===========================================================================
Sensitivity is far more interpretable in HI science as a column-density limit
than in \mjybeam. The per-channel noise is converted to a brightness
temperature via the Rayleigh--Jeans, Gaussian-beam relation
\begin{equation}
  T_{\mathrm B} = 1.222\times10^{3}\;
  \frac{S\,[\si{mJy.beam^{-1}}]}{\nu^2[\si{GHz}]\;\theta_{\mathrm{maj}}["]\,\theta_{\mathrm{min}}["]}\ \ [\si{K}],
  \label{eq:tb}
\end{equation}
and then to an HI column density through
$N_{\mathrm{HI}} = 1.823\times10^{18}\!\int T_{\mathrm B}\,\mathrm{d}v$
($\si{cm^{-2}}$, optically thin). The toolkit reports the $n\sigma$ (default
$3\sigma$) column-density sensitivity both per channel,
\begin{equation}
  N_{\mathrm{HI}}^{\,n\sigma}(\text{chan}) =
  1.823\times10^{18}\; n\; T_{\mathrm B}(\sigma)\;\delta v,
\end{equation}
and integrated over a line of width $\Delta v$ (default \SI{20}{km/s}), where
adding $\Delta v/\delta v$ independent channels gives
\begin{equation}
  N_{\mathrm{HI}}^{\,n\sigma}(\Delta v) =
  1.823\times10^{18}\; n\; T_{\mathrm B}(\sigma)\;\sqrt{\delta v\,\Delta v}\,,
  \label{eq:nhi}
\end{equation}
with $\delta v$ the channel width in velocity. The per-channel sensitivity is
plotted as a spectrum, and the integrated $3\sigma$ limit appears in the summary
table, the band CSV and the survey rollup. Line width $\Delta v$ and $n$ are set
by \code{sensitivity.nhi\_dv\_kms} and \code{sensitivity.nsigma}.

% ===========================================================================
\section{Spatial diagnostic maps and continuum residuals}\label{sec:maps}
% ===========================================================================
A single number per channel hides position-dependent problems. For the
noise-reference product the reducer accumulates, in the same streaming pass, two
2-D maps:
\begin{description}[leftmargin=1.6em,style=nextline]
  \item[Noise map] the per-pixel standard deviation \emph{across channels}.
        Calibration ripples, solar interference and residual sidelobe structure
        show up as spatial patterns that the band-median RMS cannot reveal.
  \item[Continuum-persistence map] for each pixel, the fraction of channels in
        which it exceeds $\kappa\sigma$ of that channel's robust noise
        (\code{continuum.sigma}, default $\kappa=3$). Thermal noise is exceeded
        in only a few percent of channels; a pixel significant in a \emph{large}
        fraction of channels is a continuum source that survived continuum
        subtraction. The headline scalar --- the percentage of pixels persistent
        in $>50\%$ of channels --- is a direct continuum-subtraction-quality
        indicator and drives a checklist criterion (Section~\ref{sec:checklist}).
\end{description}
If a cube is primary-beam corrected (detected from \code{BUNIT}/filename), the
report flags that its noise is position-dependent and the single-number RMS is
only indicative.

% ===========================================================================
\section{Imaging provenance}\label{sec:provenance}
% ===========================================================================
\code{provenance.py} reads how each cube was made from the FITS header and
\code{HISTORY}: the Briggs \code{robust} value (\code{WSCWEIGH} keyword or the
\code{-weight briggs} history token), \code{niter}, threshold, restoring beam,
and observation metadata (\code{TELESCOP}, \code{OBJECT}, \code{DATE-OBS},
\code{BUNIT}). These are stamped on each unit so a reviewer can see the imaging
settings at a glance. When \code{radiometer.nfac: auto}, the robust value is
mapped to an approximate weighting penalty (Section~\ref{sec:radiometer}), so the
theory line reflects the actual weighting without manual bookkeeping.

% ===========================================================================
\section{Performance and memory}\label{sec:performance}
% ===========================================================================
HI cubes are large (thousands of channels, $\sim$$10^3\times10^3$ pixels;
tens of GB). Two design choices keep the toolkit fast and memory-light:
\begin{enumerate}[noitemsep,topsep=2pt]
  \item \textbf{Memory-mapped, chunked I/O.} \code{reduction.py} memory-maps each
        cube and processes it in blocks of \code{reduce.chunk} channels, so only
        a few hundred MB are resident at once regardless of cube size. The whole
        cube is never held in RAM, and the per-cube peak is one chunk plus the
        2-D accumulators.
  \item \textbf{Vectorised statistics.} the per-channel MAD-RMS, peak, mean,
        integrated flux, 2-D maps and continuum counts are computed with
        \textsc{NumPy} reductions over each block, replacing the previous
        Python-level per-channel loop. In practice this is roughly one to two
        orders of magnitude faster.
\end{enumerate}
A single streaming pass yields all the 1-D spectra and 2-D maps; the diagnostic
channel planes and the boresight box are then read individually from the memory
map, so no second full pass is needed.

\paragraph{Plot resolution (downsampled vs full).} An HI cube can have several
thousand channels, and drawing every point of every trace makes the interactive
HTML large and sluggish (tens of MB per session). By default the line plots are
therefore \emph{decimated} to at most \code{plots.max\_points} points per trace
(default $2000$): a uniform stride over channels, but with \emph{every Bad
channel kept} so outliers are never hidden. This affects only the on-screen
plots --- the CSV tables and all statistics are computed at full spectral
resolution. The user chooses between the two versions: the downsampled default,
or a full-resolution report via \code{plots.full\_resolution: true} in the config
or the \code{-{}-full-res} command-line flag (with \code{-{}-downsample} to force
the default). The active mode is logged at the start of the run.

% ===========================================================================
\section{Flagging summary}\label{sec:flags}
% ===========================================================================
\textsc{CARACal} writes per-stage flagging-summary JSON files (band-pass and
gain calibrators, static masks, per-pointing RFI, etc.). \code{hi\_cube\_qa}
globs these per session (\code{flagreport\_glob}), recursively extracts every
flagged fraction, and renders a horizontal bar chart colour-coded against
warning ($>10\%$) and severe ($>30\%$) thresholds. The same report also drives
the effective-time reduction in the per-channel theory
(Eq.~\ref{eq:theoryflag}, Section~\ref{sec:flagaware}), so this panel is shown
once at the session level alongside the per-cube noise diagnostics.

% ===========================================================================
\section{The expert checklist}\label{sec:checklist}
% ===========================================================================
For each cube product the toolkit evaluates a fixed set of pass/warn/fail
criteria (\code{build\_\allowbreak expert\_\allowbreak checklist()}), summarising the figures into an
at-a-glance verdict. Table~\ref{tab:checklist} lists the criteria and their
scientific rationale.

\begin{longtable}{@{}L{0.30\textwidth} L{0.34\textwidth} L{0.30\textwidth}@{}}
\caption{Expert QA criteria and their justification.}\label{tab:checklist}\\
\toprule
\textbf{Criterion} & \textbf{Test} & \textbf{Why it matters}\\
\midrule
\endfirsthead
\toprule
\textbf{Criterion} & \textbf{Test} & \textbf{Why it matters}\\
\midrule
\endhead
\bottomrule
\endfoot
Noise within $1.5\times$ theory &
$R < 1.5$ pass, $<3$ warn, else fail &
The achieved sensitivity should approach the radiometer limit; a large excess
indicates calibration loss, RFI, or wrong weighting assumptions.\\[2pt]
Few bad channels ($<5\%$) &
$f_{\mathrm{bad}} < 5\%$ pass, $<15\%$ warn &
A small fraction of unusable channels is normal (band edges, persistent RFI);
a large fraction erodes the usable bandwidth.\\[2pt]
No catastrophic loss ($<15\%$ bad) &
$f_{\mathrm{bad}} < 15\%$ pass, else fail &
Beyond $\sim15\%$ the cube is compromised for line science.\\[2pt]
Noise-floor convergence &
median RMS of the last $10\%$ of channels vs.\ the first $90\%$:
$<1.2$ pass, $<2$ warn &
The noise should be stable across the band; a rising floor at the band edge
flags bandpass roll-off or edge RFI.\\[2pt]
Beam stability &
$\operatorname{std}(\theta_{\mathrm{maj}}) < \SI{0.5}{arcsec}$ pass &
A strongly varying beam prevents a common restoring beam and biases moment
maps.\\[2pt]
Reasonable axial ratio &
$\theta_{\mathrm{maj}}/\theta_{\mathrm{min}} < 2$ pass, $<3$ warn &
A highly elongated beam signals poor $uv$-coverage or a low-elevation track.\\[2pt]
No severe RMS outliers &
no channel $> 5\times$ the median RMS (0 pass, $<5$ warn) &
Isolated high-RMS channels are usually unflagged RFI spikes.\\[2pt]
Low continuum residual &
$<0.5\%$ of pixels significant in $>50\%$ of channels pass, $<2\%$ warn &
Pixels bright in most channels are continuum sources surviving continuum
subtraction (Section~\ref{sec:maps}).\\[2pt]
Gaussian noise (kurtosis) &
$<5\%$ of channels $>5\sigma_{\kappa}$ from 3 pass, $<15\%$ warn &
Non-Gaussian channels (residual RFI, incomplete continuum subtraction) that the
RMS check alone cannot see (Section~\ref{sec:kurtosis}).\\[2pt]
Flat baseline &
fitted trend of the per-channel mean, in standard errors of the mean:
$<5$ pass, $<10$ warn &
A residual slope or curvature left by continuum subtraction does not show in the noise
and appears as a seam where pointings are mosaicked.\\[2pt]
Low large-scale noise &
median ratio of the smoothed-map rms to the thermal estimate from neighbouring
channels: $<1.5$ pass, $<3$ warn &
Field-wide ripples that are present in every channel average out in the baseline and add
little to the pixel noise, but stand out in maps smoothed to a few arcminutes.\\
\end{longtable}

% ===========================================================================
\section{Configuration reference}\label{sec:config}
% ===========================================================================
A run is driven by a single YAML file. The toolkit ships a fully commented
template, regenerated with \code{hi-cube-qa -{}-example-config}. A
representative configuration:

\begin{lstlisting}
output:
    dir: ./qa_output        # one <obsid>_*.html / .csv set per session
    html: true
    csv:  true
    log:  true

sessions:                   # each root -> one HTML report
    - root: /data/G23/1610268336_sdp_l0/caracal
    - root: /data/G23/1609151470_sdp_l0/caracal

products: [dirty, image, residual]   # default three; or 'all'; or any subset
pointings: null             # null = all; or e.g. [p63, p64]
cubes: null                 # null/'all' = every cube_N; or e.g. [cube_1] / [1]

flagreport_glob: output_hi_1gc/diagnostic_plots/*-flagging-summary-*.json

radiometer:
    sefd:    425.0          # Jy per antenna, flat; if null, derived from tsys/eta
    tsys:    25.0           # K  (MeerKAT L-band typical)
    eta:     0.80           # aperture efficiency
    n_ant:          auto    # real per-session count from obsinfo; or a fixed number
    n_ant_fallback: 60      # used only if 'auto' finds nothing
    t_total: auto           # on-source hours: from obsinfo per pointing, or a number
    df:      26367.19       # Hz; auto-overridden per cube from the header
    nfac:    1.7            # weighting penalty, or 'auto' (from header robust)
    sefd_curve:    default  # 'default' freq curve | 'flat' | [[MHz,Jy],...]
    flag_fraction: auto     # fold flagging into theory: 'auto' | number | null

spectral:
    z_source: null          # e.g. 0.05 to mark a target redshift

sensitivity:
    nhi_dv_kms: 20.0        # quote 3-sigma N_HI over this line width
    nsigma:     3.0

continuum:
    sigma: 3.0              # per-channel threshold for the persistence map

reduce:
    chunk: 64               # channels read per block (memory/speed)
    maps:  true             # render 2-D noise + continuum maps

thumbnails:
    selectors: [rms, snr]   # diagnostic-channel selectors

boresight:
    enable: true
    box_size: 100           # pixels, centred on CRPIX
    centre_pix: null        # null -> use header pointing centre
    sigma_clip: 3.0
    clip_iterations: 3

thresholds:
    rms_warn_factor:   1.5
    rms_bad_factor:    3.0
    bad_chan_warn_pct: 5.0
    bad_chan_bad_pct:  15.0
    beam_scatter_warn: 0.5
    flag_warn_pct:     10.0
    flag_bad_pct:      30.0
    kurtosis_warn_sigma: 5.0   # |kurtosis-3| > N standard errors -> Warn
    kurtosis_bad_sigma:  10.0  # |kurtosis-3| > N standard errors -> Bad
\end{lstlisting}

\noindent\textbf{Important:} \code{radiometer.t\_total} is the \emph{on-source}
integration time and directly scales $\sigma_{\mathrm{theory}}$ via
Eqs.~(\ref{eq:radiometer},~\ref{eq:theoryflag}); it must be set to the real
track length for the theory comparison to be meaningful. The channel width
\code{df} is auto-read from each cube and need not be exact in the config.

% ===========================================================================
\section{Outputs}\label{sec:outputs}
% ===========================================================================
Written into \code{output.dir}:
\begin{description}[leftmargin=1.6em,style=nextline]
  \item[\code{qa\_index.html}] \textbf{The master entry point.} A small shell
        page with a session sidebar; selecting a session loads its report on
        demand in a frame, so the master stays a few kilobytes and remains
        responsive however many sessions are present. The per-session reports
        must remain in the same directory (they are referenced by relative
        name). Open this file to browse everything.
  \item[\code{<obsid>\_qa\_report.html}] One interactive report per session: a
        pointing navigation bar, a session overview table, one block per
        \code{(pointing, cube\_N)} unit (provenance strip, summary table with
        $N_{\mathrm{HI}}$, RMS curves with the per-channel theory curve,
        RMS/theory ratio, cumulative RMS, RMS histogram, a noise-overview panel
        (boresight/median/peak against the flat and flagging-reduced theory
        curves), the per-channel kurtosis panel (Section~\ref{sec:kurtosis}),
        spectral profile, integrated band spectrum, $N_{\mathrm{HI}}$
        sensitivity spectrum, beam statistics, the 2-D noise and continuum
        maps, diagnostic channel maps, and the per-unit checklist), and a
        session-level flagging chart. Plotly
        \citep{plotly} loads once from CDN; all images are embedded as base-64
        PNGs, so the file is self-contained.
  \item[\code{<obsid>\_channel\_qa.csv}] One row per channel per product:
        frequency, velocity, RMS, RMS/theory, quality, kurtosis, and beam
        parameters.
  \item[\code{<obsid>\_band\_summary.csv}] One row per unit/product with the
        band-level noise, $N_{\mathrm{HI}}$ sensitivity, continuum-persistence,
        kurtosis (\code{pct\_bad\_kurtosis}, \code{pct\_warn\_kurtosis}) and beam
        statistics.
  \item[\code{<obsid>\_checklist.csv}] The expert-checklist verdicts.
  \item[\code{qa\_survey\_summary.\{json,csv\}}] A \textbf{survey-level rollup}:
        one row per unit/product across \emph{all} sessions (median RMS,
        RMS/theory, $N_{\mathrm{HI}}$, continuum, PB-correction flag, grade), for
        tracking quality over the whole campaign.
  \item[\code{pngs/<obsid>/<pointing>\_<cube>\_<figure>.png}] When
        \code{output.save\_png} is set, the key figures (configurable via
        \code{output.png\_figures}) are also written as standalone PNGs for
        slides and logbooks. This needs the \code{kaleido} package with a
        working Chromium underneath it (both can be true even when
        \code{kaleido} is \emph{installed}, e.g.\ a version mismatch with
        Plotly, or a sandboxed HPC/Singularity environment) -- a functional
        probe render is attempted once at the start of the run, and if it
        fails, PNG export is disabled for the entire run with a single clear
        log message rather than repeating the same failed attempt once per
        session; the \code{pngs/} directory is then never created at all,
        rather than being left behind empty. Figures containing named HTML
        entities (e.g.\ \code{\&sigma;}, resolved correctly by any real
        browser for the HTML report) are exported to PNG using a sanitised
        copy with the literal Unicode character substituted, since Kaleido's
        renderer does not reliably resolve them and would otherwise fail that
        figure specifically. None of this affects the interactive HTML, which
        renders independently in the browser.
  \item[\code{hi\_cube\_qa.log}] The full run log.
\end{description}

% ===========================================================================
\section{Running the toolkit}\label{sec:running}
% ===========================================================================
\paragraph{Preview (no heavy compute).} List exactly which units and products
will be assessed:
\begin{lstlisting}
hi-cube-qa --config qa_G23_all_sessions.yaml --list
\end{lstlisting}
\paragraph{Full run, locally.}
\begin{lstlisting}
hi-cube-qa --config qa_G23_all_sessions.yaml
hi-cube-qa --config qa_G23_all_sessions.yaml --full-res   # every channel
\end{lstlisting}
\paragraph{On the cluster (Slurm + Singularity).}
\begin{lstlisting}
sbatch run_qa.sh                  # uses qa_G23_all_sessions.yaml
sbatch run_qa.sh my_config.yaml   # or a custom config
\end{lstlisting}
Cubes are loaded one unit at a time, so peak memory is a few cubes rather than a
whole session. The toolkit is built on \code{numpy} \citep{harris2020},
\code{scipy} \citep{virtanen2020}, \code{astropy} \citep{astropy2022},
\code{matplotlib} \citep{hunter2007} and \code{plotly} \citep{plotly} (all
present in the ILIFU \code{ASTRO-PY3.12} container).

% ===========================================================================
\appendix
\section{Symbol glossary}\label{app:symbols}
% ===========================================================================
\begin{longtable}{@{}l L{0.74\textwidth}@{}}
\toprule
\textbf{Symbol} & \textbf{Meaning}\\
\midrule
\endhead
$\sigma_{\mathrm{MAD}}$ & Robust per-channel RMS noise, Eq.~(\ref{eq:mad}).\\
$\sigma_{\mathrm{theory}}(\nu)$ & Per-channel radiometer noise, frequency- and flagging-aware, Eq.~(\ref{eq:theoryflag}).\\
$R$ & Band-averaged RMS-to-theory ratio (median \textsf{Good}-channel RMS over $\sigma_{\mathrm{theory}}$).\\
$\mathrm{SEFD}(\nu)$ & System-equivalent flux density vs frequency, Eq.~(\ref{eq:sefd}) and Section~\ref{sec:sefd}.\\
$f_{\mathrm{flag}}(\nu)$, $\tau_{\mathrm{eff}}$ & Flagged fraction and flagging-reduced effective time, Eq.~(\ref{eq:theoryflag}).\\
$T_{\mathrm B}$, $N_{\mathrm{HI}}$ & Brightness temperature and HI column density, Eqs.~(\ref{eq:tb},~\ref{eq:nhi}).\\
$T_{\mathrm{sys}}$, $\eta$ & System temperature and aperture efficiency.\\
$d$, $A_{\mathrm{geo}}$ & Dish diameter ($\SI{13.5}{m}$) and geometric area.\\
$N_{\mathrm{ant}}$, $n_{\mathrm{pair}}$ & Antenna count and number of ordered antenna pairs.\\
$\tau$, $\Delta\nu$, $n_{\mathrm{pol}}$ & On-source time, channel width, polarisations ($2$).\\
$n_{\mathrm{fac}}$ & Imaging-weight noise penalty relative to natural.\\
$\theta_{\mathrm{maj}}$, $\theta_{\mathrm{min}}$ & Beam major/minor FWHM (arcsec).\\
$\Omega_{\mathrm{beam}}$, $\Omega_{\mathrm{pix}}$, $n_{\mathrm{pb}}$ & Beam/pixel solid angle and pixels-per-beam.\\
$S_{\mathrm{int}}(\nu)$ & Spatially integrated flux-density spectrum, Eq.~(\ref{eq:sint}).\\
$f_0$, $c$ & \textsc{Hi}\ rest frequency, speed of light.\\
$f_{\mathrm{bad}}$, $f_{\mathrm{warn}}$ & Bad/warn channel fractions.\\
\bottomrule
\end{longtable}

\noindent\emph{Note on the boresight method.} The best-case boresight noise
methodology (Section~\ref{sec:boresight}) follows the reference MeerKAT
\textsc{Hi}\ reduction (referred to as ``P20'' in the source comments); replace
this note with the appropriate bibliographic reference for a formal write-up.

% ===========================================================================
\bibliographystyle{plainnat}
\bibliography{refs}
% ===========================================================================

\end{document}
